% arara: xelatex: { options: ['-synctex=1', '-interaction=nonstopmode'] }

\documentclass{standalone}
\usepackage{tikz}
\usetikzlibrary{shapes, arrows.meta, positioning, fit}
\usepackage{xeCJK}

\begin{document}
\begin{tikzpicture}[
  font=\small,
  node distance=8mm and 12mm,
  process/.style={rectangle, draw, minimum height=10mm, minimum width=32mm, text centered, fill=green!5},
  dashedlink/.style={->, thick, dashed, red},
  solidlink/.style={->, thick},
  dashedbox/.style={draw=black, dashed, inner sep=5mm, fit=#1}
]

% 第一层：输入
\node[process] (inputCode) {输入代码};
\node[process, right=of inputCode] (inputContext) {输入上下文};

% 第二层：代码切片
\node[process, below=of inputCode, xshift=15mm] (codeSlice) {代码切片};

% 第三层：代码特征识别
\node[process, below=of codeSlice] (codeFeature) {代码特征识别};

% 第四层：匹配与分析
\node[process, below left=of codeFeature, xshift=-15mm] (templateMatch) {模板匹配};
\node[process, below right=of codeFeature, xshift=15mm] (ragResult) {RAG 结果};
\node[process, below=of ragResult] (programAnalysis) {程序分析信息增强};

% 第五层：模板处理 与 模型推理 并列
\node[process, below=of programAnalysis, xshift=-40mm] (problemTemplate) {代码问题识别模板};
\node[process, below=of problemTemplate] (opinionTemplate) {检测意见生成模板};
\node[process, below=of opinionTemplate] (reflection) {检测意见生成反思};

\node[process, right=6cm of opinionTemplate] (llm) {大模型 (Owen/Codellama)};

% 第六层：输出
\node[process, below left=of llm, xshift=-20mm] (codeOutput) {代码输出结果};
\node[process, below right=of llm, xshift=20mm] (problemResult) {代码问题识别结果};

% 第七层：汇总与数据集
\node[process, below=of llm, yshift=-45mm] (resultSummary) {代码推理结果汇总};
\node[process, below=of resultSummary] (dataset) {数据集构建与效果验证};

% 原有连接
\draw[solidlink] (inputCode) -- (codeSlice);
\draw[solidlink] (inputContext) -- (codeSlice);
\draw[solidlink] (codeSlice) -- (codeFeature);
\draw[solidlink] (codeFeature) -- (templateMatch);
\draw[solidlink] (codeFeature) -- (ragResult);
\draw[solidlink] (templateMatch) -- (ragResult);
\draw[solidlink] (ragResult) -- (programAnalysis);
\draw[solidlink] (programAnalysis) -- (problemTemplate);
\draw[solidlink] (problemTemplate) -- (opinionTemplate);
\draw[solidlink] (opinionTemplate) -- (reflection);
\draw[solidlink] (reflection) -- (opinionTemplate);
\draw[solidlink] (opinionTemplate) -- (llm);
\draw[solidlink] (llm) -- (codeOutput);
\draw[solidlink] (llm) -- (problemResult);
\draw[solidlink] (codeOutput) -- (resultSummary);
\draw[solidlink] (problemResult) -- (resultSummary);
\draw[solidlink] (resultSummary) -- (dataset);

% 新增连接
\draw[dashedlink] (programAnalysis) -- (opinionTemplate);
\draw[dashedlink] (codeFeature) -- (problemTemplate);
\draw[dashedlink] (ragResult) -- (llm);
\draw[dashedlink] (resultSummary) -- (reflection);
\draw[dashedlink] (dataset) to[out=180, in=270] (problemTemplate);

% 包围框（虚线）
\node[dashedbox={(inputCode)(inputContext)(dataset)(problemResult)}] {};

% 图例
\node[draw, rectangle, fill=white, below=10mm of dataset, align=left] (legend) {
    \textbf{图例:} \\
    \tikz\draw[solidlink] (0,0) -- (0.7,0); 原有连接 \\
    \tikz\draw[dashedlink] (0,0) -- (0.7,0); 新增连接
};

\end{tikzpicture}
\end{document}